\documentclass[
  doc,
  floatsintext,
  longtable,
  a4paper,
  nolmodern,
  notxfonts,
  notimes,
  12pt,
  colorlinks=true,linkcolor=blue,citecolor=blue,urlcolor=blue]{apa7}

\usepackage{amsmath}
\usepackage{amssymb}

\geometry{inner=1in, outer=1in}
\fancyhfoffset[LE,RO]{0cm}


\usepackage[bidi=default]{babel}
\babelprovide[main,import]{spanish}


% get rid of language-specific shorthands (see #6817):
\let\LanguageShortHands\languageshorthands
\def\languageshorthands#1{}

\RequirePackage{longtable}
\RequirePackage{threeparttablex}

\makeatletter
\renewcommand{\paragraph}{\@startsection{paragraph}{4}{\parindent}%
	{0\baselineskip \@plus 0.2ex \@minus 0.2ex}%
	{-.5em}%
	{\normalfont\normalsize\bfseries\typesectitle}}

\renewcommand{\subparagraph}[1]{\@startsection{subparagraph}{5}{0.5em}%
	{0\baselineskip \@plus 0.2ex \@minus 0.2ex}%
	{-\z@\relax}%
	{\normalfont\normalsize\bfseries\itshape\hspace{\parindent}{#1}\textit{\addperi}}{\relax}}
\makeatother




\usepackage{longtable, booktabs, multirow, multicol, colortbl, hhline, caption, array, float, xpatch}
\usepackage{subcaption}


\renewcommand\thesubfigure{\Alph{subfigure}}
\setcounter{topnumber}{2}
\setcounter{bottomnumber}{2}
\setcounter{totalnumber}{4}
\renewcommand{\topfraction}{0.85}
\renewcommand{\bottomfraction}{0.85}
\renewcommand{\textfraction}{0.15}
\renewcommand{\floatpagefraction}{0.7}

\usepackage{tcolorbox}
\tcbuselibrary{listings,theorems, breakable, skins}
\usepackage{fontawesome5}

\definecolor{quarto-callout-color}{HTML}{909090}
\definecolor{quarto-callout-note-color}{HTML}{0758E5}
\definecolor{quarto-callout-important-color}{HTML}{CC1914}
\definecolor{quarto-callout-warning-color}{HTML}{EB9113}
\definecolor{quarto-callout-tip-color}{HTML}{00A047}
\definecolor{quarto-callout-caution-color}{HTML}{FC5300}
\definecolor{quarto-callout-color-frame}{HTML}{ACACAC}
\definecolor{quarto-callout-note-color-frame}{HTML}{4582EC}
\definecolor{quarto-callout-important-color-frame}{HTML}{D9534F}
\definecolor{quarto-callout-warning-color-frame}{HTML}{F0AD4E}
\definecolor{quarto-callout-tip-color-frame}{HTML}{02B875}
\definecolor{quarto-callout-caution-color-frame}{HTML}{FD7E14}

%\newlength\Oldarrayrulewidth
%\newlength\Oldtabcolsep


\usepackage{hyperref}



\usepackage{color}
\usepackage{fancyvrb}
\newcommand{\VerbBar}{|}
\newcommand{\VERB}{\Verb[commandchars=\\\{\}]}
\DefineVerbatimEnvironment{Highlighting}{Verbatim}{commandchars=\\\{\}}
% Add ',fontsize=\small' for more characters per line
\usepackage{framed}
\definecolor{shadecolor}{RGB}{241,243,245}
\newenvironment{Shaded}{\begin{snugshade}}{\end{snugshade}}
\newcommand{\AlertTok}[1]{\textcolor[rgb]{0.68,0.00,0.00}{#1}}
\newcommand{\AnnotationTok}[1]{\textcolor[rgb]{0.37,0.37,0.37}{#1}}
\newcommand{\AttributeTok}[1]{\textcolor[rgb]{0.40,0.45,0.13}{#1}}
\newcommand{\BaseNTok}[1]{\textcolor[rgb]{0.68,0.00,0.00}{#1}}
\newcommand{\BuiltInTok}[1]{\textcolor[rgb]{0.00,0.23,0.31}{#1}}
\newcommand{\CharTok}[1]{\textcolor[rgb]{0.13,0.47,0.30}{#1}}
\newcommand{\CommentTok}[1]{\textcolor[rgb]{0.37,0.37,0.37}{#1}}
\newcommand{\CommentVarTok}[1]{\textcolor[rgb]{0.37,0.37,0.37}{\textit{#1}}}
\newcommand{\ConstantTok}[1]{\textcolor[rgb]{0.56,0.35,0.01}{#1}}
\newcommand{\ControlFlowTok}[1]{\textcolor[rgb]{0.00,0.23,0.31}{\textbf{#1}}}
\newcommand{\DataTypeTok}[1]{\textcolor[rgb]{0.68,0.00,0.00}{#1}}
\newcommand{\DecValTok}[1]{\textcolor[rgb]{0.68,0.00,0.00}{#1}}
\newcommand{\DocumentationTok}[1]{\textcolor[rgb]{0.37,0.37,0.37}{\textit{#1}}}
\newcommand{\ErrorTok}[1]{\textcolor[rgb]{0.68,0.00,0.00}{#1}}
\newcommand{\ExtensionTok}[1]{\textcolor[rgb]{0.00,0.23,0.31}{#1}}
\newcommand{\FloatTok}[1]{\textcolor[rgb]{0.68,0.00,0.00}{#1}}
\newcommand{\FunctionTok}[1]{\textcolor[rgb]{0.28,0.35,0.67}{#1}}
\newcommand{\ImportTok}[1]{\textcolor[rgb]{0.00,0.46,0.62}{#1}}
\newcommand{\InformationTok}[1]{\textcolor[rgb]{0.37,0.37,0.37}{#1}}
\newcommand{\KeywordTok}[1]{\textcolor[rgb]{0.00,0.23,0.31}{\textbf{#1}}}
\newcommand{\NormalTok}[1]{\textcolor[rgb]{0.00,0.23,0.31}{#1}}
\newcommand{\OperatorTok}[1]{\textcolor[rgb]{0.37,0.37,0.37}{#1}}
\newcommand{\OtherTok}[1]{\textcolor[rgb]{0.00,0.23,0.31}{#1}}
\newcommand{\PreprocessorTok}[1]{\textcolor[rgb]{0.68,0.00,0.00}{#1}}
\newcommand{\RegionMarkerTok}[1]{\textcolor[rgb]{0.00,0.23,0.31}{#1}}
\newcommand{\SpecialCharTok}[1]{\textcolor[rgb]{0.37,0.37,0.37}{#1}}
\newcommand{\SpecialStringTok}[1]{\textcolor[rgb]{0.13,0.47,0.30}{#1}}
\newcommand{\StringTok}[1]{\textcolor[rgb]{0.13,0.47,0.30}{#1}}
\newcommand{\VariableTok}[1]{\textcolor[rgb]{0.07,0.07,0.07}{#1}}
\newcommand{\VerbatimStringTok}[1]{\textcolor[rgb]{0.13,0.47,0.30}{#1}}
\newcommand{\WarningTok}[1]{\textcolor[rgb]{0.37,0.37,0.37}{\textit{#1}}}

\providecommand{\tightlist}{%
  \setlength{\itemsep}{0pt}\setlength{\parskip}{0pt}}
\usepackage{longtable,booktabs,array}
\newcounter{none} % for unnumbered tables
\usepackage{calc} % for calculating minipage widths
% Correct order of tables after \paragraph or \subparagraph
\usepackage{etoolbox}
\makeatletter
\patchcmd\longtable{\par}{\if@noskipsec\mbox{}\fi\par}{}{}
\makeatother
% Allow footnotes in longtable head/foot
\IfFileExists{footnotehyper.sty}{\usepackage{footnotehyper}}{\usepackage{footnote}}
\makesavenoteenv{longtable}

\usepackage{graphicx}
\makeatletter
\newsavebox\pandoc@box
\newcommand*\pandocbounded[1]{% scales image to fit in text height/width
  \sbox\pandoc@box{#1}%
  \Gscale@div\@tempa{\textheight}{\dimexpr\ht\pandoc@box+\dp\pandoc@box\relax}%
  \Gscale@div\@tempb{\linewidth}{\wd\pandoc@box}%
  \ifdim\@tempb\p@<\@tempa\p@\let\@tempa\@tempb\fi% select the smaller of both
  \ifdim\@tempa\p@<\p@\scalebox{\@tempa}{\usebox\pandoc@box}%
  \else\usebox{\pandoc@box}%
  \fi%
}
% Set default figure placement to htbp
\def\fps@figure{htbp}
\makeatother







\usepackage{newtx}

\defaultfontfeatures{Scale=MatchLowercase}
\defaultfontfeatures[\rmfamily]{Ligatures=TeX,Scale=1}





\title{Sintaxis de Python}


\shorttitle{SINTAXIS DE PYTHON}


\usepackage{etoolbox}



\ccoppy{\textcopyright~2026}



\author{Edison Achalma}



\affiliation{
{Escuela Profesional de Economía, Universidad Nacional de San Cristóbal
de Huamanga}}




\leftheader{Achalma}

\date{2026-10-20}


\abstract{Este abstract será actualizado una vez que se complete el
contenido final del artículo. }



\authornote{\par{\addORCIDlink{Edison Achalma}{0000-0001-6996-3364}} 
\par{ }
\par{   El autor no tiene conflictos de interés que revelar.    Los
roles de autor se clasificaron utilizando la taxonomía de roles de
colaborador (CRediT; https://credit.niso.org/) de la siguiente
manera:  Edison Achalma:   conceptualización, metodología, análisis
formal, investigación, recursos, curación de
datos, redacción, visualización, supervisión, administración del
proyecto}
\par{La correspondencia relativa a este artículo debe dirigirse a Edison
Achalma, Escuela Profesional de Economía, Universidad Nacional de San
Cristóbal de Huamanga, Portal Independencia N°
57, Ayacucho, AYA 5001, Perú, Email: \href{mailto:elmer.achalma.09@unsch.edu.pe}{elmer.achalma.09@unsch.edu.pe}}
}

\makeatletter
\let\endoldlt\endlongtable
\def\endlongtable{
\hline
\endoldlt
}
\makeatother

\urlstyle{same}



\hypersetup{
  pdftitle={Sintaxis de Python},
  pdfauthor={Autor}
}
\makeatletter
\@ifpackageloaded{caption}{}{\usepackage{caption}}
\AtBeginDocument{%
\ifdefined\contentsname
  \renewcommand*\contentsname{Tabla de contenidos}
\else
  \newcommand\contentsname{Tabla de contenidos}
\fi
\ifdefined\listfigurename
  \renewcommand*\listfigurename{Lista de Figuras}
\else
  \newcommand\listfigurename{Lista de Figuras}
\fi
\ifdefined\listtablename
  \renewcommand*\listtablename{Lista de Tablas}
\else
  \newcommand\listtablename{Lista de Tablas}
\fi
\ifdefined\figurename
  \renewcommand*\figurename{Figura}
\else
  \newcommand\figurename{Figura}
\fi
\ifdefined\tablename
  \renewcommand*\tablename{Tabla}
\else
  \newcommand\tablename{Tabla}
\fi
}
\@ifpackageloaded{float}{}{\usepackage{float}}
\floatstyle{ruled}
\@ifundefined{c@chapter}{\newfloat{codelisting}{h}{lop}}{\newfloat{codelisting}{h}{lop}[chapter]}
\floatname{codelisting}{Listado}
\newcommand*\listoflistings{\listof{codelisting}{Lista de Listados}}
\makeatother
\makeatletter
\makeatother
\makeatletter
\@ifpackageloaded{caption}{}{\usepackage{caption}}
\@ifpackageloaded{subcaption}{}{\usepackage{subcaption}}
\makeatother
\makeatletter
\@ifpackageloaded{fontawesome5}{}{\usepackage{fontawesome5}}
\makeatother

% From https://tex.stackexchange.com/a/645996/211326
%%% apa7 doesn't want to add appendix section titles in the toc
%%% let's make it do it
\makeatletter
\xpatchcmd{\appendix}
  {\par}
  {\addcontentsline{toc}{section}{\@currentlabelname}\par}
  {}{}
\makeatother

%% Disable longtable counter
%% https://tex.stackexchange.com/a/248395/211326

\usepackage{etoolbox}

\makeatletter
\patchcmd{\LT@caption}
  {\bgroup}
  {\bgroup\global\LTpatch@captiontrue}
  {}{}
\patchcmd{\longtable}
  {\par}
  {\par\global\LTpatch@captionfalse}
  {}{}
\apptocmd{\endlongtable}
  {\ifLTpatch@caption\else\addtocounter{table}{-1}\fi}
  {}{}
\newif\ifLTpatch@caption
\makeatother

\begin{document}

\maketitle


\hypertarget{toc}{}
\tableofcontents
\newpage
\section[Introduction]{Sintaxis de Python}

\setcounter{secnumdepth}{3}

\setlength\LTleft{0pt}




\subsection{🎯 Motivación}\label{motivaciuxf3n}

La sintaxis es el conjunto de reglas que definen cómo se escribe código
válido en Python. A diferencia de otros lenguajes donde la sintaxis es
más permisiva, Python fue diseñado para ser \textbf{legible y
consistente}. Entender la sintaxis desde el principio te ahorrará
errores frustrantes y hará que tu código sea más claro para otros (y
para ti mismo cuando lo releas).

\subsection{💡 Intuición}\label{intuiciuxf3n}

Piensa en la sintaxis de Python como las reglas gramaticales de un
idioma. En español, ``El perro come'' es válido; ``Perro el come'' suena
mal. En Python, la indentación es como la puntuación: no es opcional, es
parte de la gramática. Si olvidas los espacios, Python no entiende dónde
termina un bloque y empieza otro, igual que si escribieras una oración
sin puntos.

\subsection{🧠 Modelo mental}\label{modelo-mental}

La sintaxis de Python se basa en tres pilares fundamentales:

\begin{enumerate}
\def\labelenumi{\arabic{enumi}.}
\item
  \textbf{Indentación como estructura}: en lugar de llaves
  (\texttt{\{\}}) o palabras clave (\texttt{begin/end}), Python usa
  espacios para agrupar código. Esto obliga a escribir código limpio y
  legible.
\item
  \textbf{Instrucciones línea por línea}: cada línea suele ser una
  instrucción completa. No necesitas punto y coma (\texttt{;}) al final.
\item
  \textbf{Comentarios documentan, no ejecutan}: todo lo que escribas
  después de \texttt{\#} es ignorado por Python. Son para ti y para
  otros humanos.
\end{enumerate}

\subsection{📐 Desarrollo}\label{desarrollo}

\subsubsection{1. Indentación: la regla de
oro}\label{indentaciuxf3n-la-regla-de-oro}

En Python, la indentación \textbf{define bloques de código}. Esto aplica
a:

\begin{itemize}
\tightlist
\item
  Condicionales (\texttt{if}, \texttt{elif}, \texttt{else})
\item
  Bucles (\texttt{for}, \texttt{while})
\item
  Funciones (\texttt{def})
\item
  Clases (\texttt{class})
\end{itemize}

\begin{Shaded}
\begin{Highlighting}[]
\ControlFlowTok{if} \DecValTok{5} \OperatorTok{\textgreater{}} \DecValTok{2}\NormalTok{:}
    \BuiltInTok{print}\NormalTok{(}\StringTok{"Cinco es mayor que dos"}\NormalTok{)  }\CommentTok{\# Esta línea está indentada}
    \BuiltInTok{print}\NormalTok{(}\StringTok{"Esto también pertenece al if"}\NormalTok{)}
\BuiltInTok{print}\NormalTok{(}\StringTok{"Esto ya está fuera del if"}\NormalTok{)   }\CommentTok{\# Sin indentación = fuera del bloque}
\end{Highlighting}
\end{Shaded}

\textbf{Reglas de indentación:}

\begin{itemize}
\tightlist
\item
  Debe haber \textbf{al menos un espacio} o un tabulador.
\item
  El \textbf{número de espacios debe ser consistente} dentro del mismo
  bloque.
\item
  La convención más común es \textbf{4 espacios} (recomendado por PEP
  8).
\end{itemize}

\subsubsection{2. Variables: creación sin
declaración}\label{variables-creaciuxf3n-sin-declaraciuxf3n}

En Python, \textbf{no declaras} variables. Se crean cuando les asignas
un valor:

\begin{Shaded}
\begin{Highlighting}[]
\NormalTok{nombre }\OperatorTok{=} \StringTok{"Ana"}      \CommentTok{\# Variable de texto (str)}
\NormalTok{edad }\OperatorTok{=} \DecValTok{25}           \CommentTok{\# Variable numérica (int)}
\NormalTok{precio }\OperatorTok{=} \FloatTok{99.99}      \CommentTok{\# Variable decimal (float)}
\NormalTok{es\_valido }\OperatorTok{=} \VariableTok{True}    \CommentTok{\# Variable booleana (bool)}
\end{Highlighting}
\end{Shaded}

Python infiere el tipo automáticamente. Las variables pueden cambiar de
tipo durante la ejecución (tipado dinámico).

\subsubsection{3. Comentarios: documenta tu
código}\label{comentarios-documenta-tu-cuxf3digo}

Los comentarios son líneas que Python ignora completamente. Sirven para
explicar qué hace tu código:

\begin{Shaded}
\begin{Highlighting}[]
\CommentTok{\# Esto es un comentario de una línea}
\BuiltInTok{print}\NormalTok{(}\StringTok{"Hola"}\NormalTok{)  }\CommentTok{\# También puedes poner comentarios al final de una línea}

\CommentTok{\# Los comentarios son útiles para:}
\CommentTok{\# {-} Explicar lógica compleja}
\CommentTok{\# {-} Desactivar temporalmente código}
\CommentTok{\# {-} Dejar notas para ti mismo o para otros}
\end{Highlighting}
\end{Shaded}

\subsubsection{4. Convenciones de sintaxis (PEP
8)}\label{convenciones-de-sintaxis-pep-8}

\begin{itemize}
\tightlist
\item
  \textbf{Sensibilidad a mayúsculas}: \texttt{miVariable} y
  \texttt{mivariable} son diferentes.
\item
  \textbf{Nombres en minúsculas con guiones bajos}:
  \texttt{nombre\_usuario} (no \texttt{nombreUsuario}).
\item
  \textbf{Constantes en mayúsculas}: \texttt{MAXIMO\ =\ 100}.
\item
  \textbf{Una línea en blanco} entre funciones y clases para
  legibilidad.
\end{itemize}

\subsection{🔬 Ejemplos}\label{ejemplos}

\textbf{Mínimo ---} Un programa que combina sintaxis válida:

\begin{Shaded}
\begin{Highlighting}[]
\CommentTok{\# Calcular descuento simple}
\NormalTok{precio }\OperatorTok{=} \DecValTok{100}
\NormalTok{descuento }\OperatorTok{=} \DecValTok{20}
\NormalTok{precio\_final }\OperatorTok{=}\NormalTok{ precio }\OperatorTok{{-}}\NormalTok{ descuento}
\BuiltInTok{print}\NormalTok{(}\StringTok{"Precio final:"}\NormalTok{, precio\_final)}
\CommentTok{\# Salida: Precio final: 80}
\end{Highlighting}
\end{Shaded}

\textbf{Progresivo ---} Indentación y bloques correctos vs incorrectos:

\begin{Shaded}
\begin{Highlighting}[]
\CommentTok{\# ✅ CORRECTO {-} mismo número de espacios (4)}
\ControlFlowTok{if} \DecValTok{10} \OperatorTok{\textgreater{}} \DecValTok{5}\NormalTok{:}
    \BuiltInTok{print}\NormalTok{(}\StringTok{"Diez es mayor"}\NormalTok{)}
    \BuiltInTok{print}\NormalTok{(}\StringTok{"Este mensaje también está dentro del if"}\NormalTok{)}

\CommentTok{\# ❌ INCORRECTO {-} indentación inconsistente}
\ControlFlowTok{if} \DecValTok{10} \OperatorTok{\textgreater{}} \DecValTok{5}\NormalTok{:}
    \BuiltInTok{print}\NormalTok{(}\StringTok{"Diez es mayor"}\NormalTok{)}
  \BuiltInTok{print}\NormalTok{(}\StringTok{"Este mensaje tiene 2 espacios"}\NormalTok{)  }\CommentTok{\# Error: indentación inconsistente}
\end{Highlighting}
\end{Shaded}

\textbf{Real ---} Un script que usa sintaxis básica con comentarios:

\begin{Shaded}
\begin{Highlighting}[]
\CommentTok{\# config.py {-} Configuración simple de usuario}
\NormalTok{nombre }\OperatorTok{=} \BuiltInTok{input}\NormalTok{(}\StringTok{"¿Cómo te llamas? "}\NormalTok{)  }\CommentTok{\# Captura entrada del usuario}
\BuiltInTok{print}\NormalTok{(}\SpecialStringTok{f"Hola, }\SpecialCharTok{\{}\NormalTok{nombre}\SpecialCharTok{\}}\SpecialStringTok{"}\NormalTok{)            }\CommentTok{\# f{-}string: formatea cadenas}

\CommentTok{\# Ejecución condicional}
\ControlFlowTok{if}\NormalTok{ nombre }\OperatorTok{==} \StringTok{"admin"}\NormalTok{:}
    \BuiltInTok{print}\NormalTok{(}\StringTok{"Modo administrador activado"}\NormalTok{)}
    \BuiltInTok{print}\NormalTok{(}\StringTok{"Tienes acceso completo"}\NormalTok{)}
\ControlFlowTok{else}\NormalTok{:}
    \BuiltInTok{print}\NormalTok{(}\StringTok{"Modo usuario"}\NormalTok{)}
    \BuiltInTok{print}\NormalTok{(}\StringTok{"Acceso limitado"}\NormalTok{)}
\end{Highlighting}
\end{Shaded}

\subsection{⚠️ Errores comunes}\label{errores-comunes}

{\def\LTcaptype{none} % do not increment counter
\begin{longtable}[]{@{}
  >{\raggedright\arraybackslash}p{(\linewidth - 4\tabcolsep) * \real{0.2184}}
  >{\raggedright\arraybackslash}p{(\linewidth - 4\tabcolsep) * \real{0.4080}}
  >{\raggedright\arraybackslash}p{(\linewidth - 4\tabcolsep) * \real{0.3736}}@{}}
\toprule\noalign{}
\begin{minipage}[b]{\linewidth}\raggedright
Error
\end{minipage} & \begin{minipage}[b]{\linewidth}\raggedright
Síntoma
\end{minipage} & \begin{minipage}[b]{\linewidth}\raggedright
Cómo detectarlo
\end{minipage} \\
\midrule\noalign{}
\endhead
\bottomrule\noalign{}
\endlastfoot
\textbf{Indentación faltante} &
\texttt{IndentationError:\ expected\ an\ indented\ block} & Busca la
línea marcada; debe tener espacios antes del código \\
\textbf{Indentación inconsistente} &
\texttt{IndentationError:\ unindent\ does\ not\ match\ any\ outer\ indentation\ level}
& Usa espacios o tabs, pero NO mezcles ambos \\
\textbf{Confundir \texttt{=} con \texttt{==}} & La variable se reasigna
en lugar de comparar & Revisa condicionales: \texttt{=} es asignación,
\texttt{==} es comparación \\
\textbf{Olvidar cerrar paréntesis/comillas} &
\texttt{SyntaxError:\ unexpected\ EOF\ while\ parsing} & Cuenta
paréntesis y comillas; deben abrirse y cerrarse \\
\textbf{Espacios en nombres de variables} &
\texttt{SyntaxError:\ invalid\ syntax} & Las variables no pueden tener
espacios: \texttt{mi\ nombre} → \texttt{mi\_nombre} \\
\end{longtable}
}

\subsection{✅ Buenas prácticas}\label{buenas-pruxe1cticas}

\begin{itemize}
\tightlist
\item
  \textbf{Usa 4 espacios para indentar} (configura tu editor para que
  TAB inserte 4 espacios).
\item
  \textbf{Comenta el ``por qué''}, no el ``qué'': el código ya dice qué
  hace.
\item
  \textbf{Nombra variables de forma descriptiva}:
  \texttt{precio\_con\_descuento} en lugar de \texttt{pcd}.
\item
  \textbf{No uses palabras reservadas} como variables (\texttt{print},
  \texttt{if}, \texttt{while}, \texttt{class}).
\item
  \textbf{Agrupa lógica relacionada} con líneas en blanco para separar
  secciones.
\end{itemize}

\subsection{📝 Ejercicio}\label{ejercicio}

\textbf{Básico} --- Corrección de errores

El siguiente código tiene \textbf{4 errores de sintaxis}. Encuéntralos y
corrígelos:

\begin{Shaded}
\begin{Highlighting}[]
\CommentTok{\# calculadora simple}
\NormalTok{numero1 }\OperatorTok{=} \DecValTok{10}
\NormalTok{numero2 }\OperatorTok{=} \DecValTok{5}
\NormalTok{suma }\OperatorTok{=}\NormalTok{ numero1 }\OperatorTok{+}\NormalTok{ numero2}
\BuiltInTok{print}\NormalTok{(}\StringTok{"La suma es"}\NormalTok{, suma)}
\ControlFlowTok{if}\NormalTok{ suma }\OperatorTok{\textgreater{}} \DecValTok{10}\NormalTok{:}
\BuiltInTok{print}\NormalTok{(}\StringTok{"La suma es mayor que 10"}\NormalTok{)}
\NormalTok{    resta }\OperatorTok{=}\NormalTok{ numero1 }\OperatorTok{{-}}\NormalTok{ numero2}
\BuiltInTok{print}\NormalTok{(}\StringTok{"La resta es"}\NormalTok{, resta)}
   \BuiltInTok{print}\NormalTok{(}\StringTok{"Fin del programa"}\NormalTok{)}
\end{Highlighting}
\end{Shaded}

\textbf{Pista:} busca problemas de indentación y nombres mal escritos.

\begin{center}\rule{0.5\linewidth}{0.5pt}\end{center}

\textbf{Solución (no mires hasta haberlo intentado):}

\begin{Shaded}
\begin{Highlighting}[]
\CommentTok{\# calculadora simple}
\NormalTok{numero1 }\OperatorTok{=} \DecValTok{10}
\NormalTok{numero2 }\OperatorTok{=} \DecValTok{5}
\NormalTok{suma }\OperatorTok{=}\NormalTok{ numero1 }\OperatorTok{+}\NormalTok{ numero2}
\BuiltInTok{print}\NormalTok{(}\StringTok{"La suma es"}\NormalTok{, suma)}
\ControlFlowTok{if}\NormalTok{ suma }\OperatorTok{\textgreater{}} \DecValTok{10}\NormalTok{:}
    \BuiltInTok{print}\NormalTok{(}\StringTok{"La suma es mayor que 10"}\NormalTok{)  }\CommentTok{\# Faltaba indentación}
\NormalTok{    resta }\OperatorTok{=}\NormalTok{ numero1 }\OperatorTok{{-}}\NormalTok{ numero2}
\BuiltInTok{print}\NormalTok{(}\StringTok{"La resta es"}\NormalTok{, resta)}
\BuiltInTok{print}\NormalTok{(}\StringTok{"Fin del programa"}\NormalTok{)  }\CommentTok{\# La indentación extra no es necesaria}
\end{Highlighting}
\end{Shaded}

\subsection{🗂️ Síntesis}\label{suxedntesis}

\begin{itemize}
\tightlist
\item
  \textbf{Indentación = bloques de código}: esencial, no opcional (4
  espacios por convención).
\item
  \textbf{Variables se crean al asignar}: no necesitan declaración
  previa.
\item
  \textbf{Comentarios con \texttt{\#}}: documentan, no ejecutan.
\item
  \textbf{Sintaxis sensible a mayúsculas}: \texttt{Variable} ≠
  \texttt{variable}.
\item
  \textbf{Usa PEP 8}: guía de estilo para código limpio y consistente.
\end{itemize}

\subsection{🔗 Qué sigue}\label{quuxe9-sigue}

\begin{itemize}
\tightlist
\item
  \textbf{Relacionado:} variables y tipos de datos (manejo de texto,
  números, booleanos)
\item
  \textbf{Siguiente tema:} Variables y tipos de datos --- aprenderás a
  almacenar y manipular información con distintos tipos de datos en
  Python.
\end{itemize}

\section{Publicaciones Similares}\label{publicaciones-similares}

Si te interesó este artículo, te recomendamos que explores otros blogs y
recursos relacionados que pueden ampliar tus conocimientos. Aquí te dejo
algunas sugerencias:

\begin{enumerate}
\def\labelenumi{\arabic{enumi}.}
\tightlist
\item
  \href{https://numerus-scriptum.netlify.app/python/2020-06-19-instalacion-de-anaconda/index.pdf}{\faIcon{file-pdf}}
  \href{https://numerus-scriptum.netlify.app/python/2020-06-19-instalacion-de-anaconda}{Instalacion
  De Anaconda}
\item
  \href{https://numerus-scriptum.netlify.app/python/2020-06-20-configurar-entorno-virtual-python-anaconda/index.pdf}{\faIcon{file-pdf}}
  \href{https://numerus-scriptum.netlify.app/python/2020-06-20-configurar-entorno-virtual-python-anaconda}{Configurar
  Entorno Virtual Python Anaconda}
\item
  \href{https://numerus-scriptum.netlify.app/python/2021-04-17-01-introducion-a-la-programacion-con-python/index.pdf}{\faIcon{file-pdf}}
  \href{https://numerus-scriptum.netlify.app/python/2021-04-17-01-introducion-a-la-programacion-con-python}{01
  Introducion A La Programacion Con Python}
\item
  \href{https://numerus-scriptum.netlify.app/python/2021-05-31-02-variables-expresiones-y-statements-con-python/index.pdf}{\faIcon{file-pdf}}
  \href{https://numerus-scriptum.netlify.app/python/2021-05-31-02-variables-expresiones-y-statements-con-python}{02
  Variables Expresiones Y Statements Con Python}
\item
  \href{https://numerus-scriptum.netlify.app/python/2021-06-07-03-objetos-de-python/index.pdf}{\faIcon{file-pdf}}
  \href{https://numerus-scriptum.netlify.app/python/2021-06-07-03-objetos-de-python}{03
  Objetos De Python}
\item
  \href{https://numerus-scriptum.netlify.app/python/2021-06-14-04-ejecucion-condicional-con-python/index.pdf}{\faIcon{file-pdf}}
  \href{https://numerus-scriptum.netlify.app/python/2021-06-14-04-ejecucion-condicional-con-python}{04
  Ejecucion Condicional Con Python}
\item
  \href{https://numerus-scriptum.netlify.app/python/2021-06-21-05-iteraciones-con-python/index.pdf}{\faIcon{file-pdf}}
  \href{https://numerus-scriptum.netlify.app/python/2021-06-21-05-iteraciones-con-python}{05
  Iteraciones Con Python}
\item
  \href{https://numerus-scriptum.netlify.app/python/2021-08-16-06-funciones-con-python/index.pdf}{\faIcon{file-pdf}}
  \href{https://numerus-scriptum.netlify.app/python/2021-08-16-06-funciones-con-python}{06
  Funciones Con Python}
\item
  \href{https://numerus-scriptum.netlify.app/python/2021-08-23-07-dataframes-con-python/index.pdf}{\faIcon{file-pdf}}
  \href{https://numerus-scriptum.netlify.app/python/2021-08-23-07-dataframes-con-python}{07
  Dataframes Con Python}
\item
  \href{https://numerus-scriptum.netlify.app/python/2021-11-29-08-prediccion-y-metrica-de-performance-con-python/index.pdf}{\faIcon{file-pdf}}
  \href{https://numerus-scriptum.netlify.app/python/2021-11-29-08-prediccion-y-metrica-de-performance-con-python}{08
  Prediccion Y Metrica De Performance Con Python}
\item
  \href{https://numerus-scriptum.netlify.app/python/2021-12-06-09-metodos-de-machine-learning-para-clasificacion-con-python/index.pdf}{\faIcon{file-pdf}}
  \href{https://numerus-scriptum.netlify.app/python/2021-12-06-09-metodos-de-machine-learning-para-clasificacion-con-python}{09
  Metodos De Machine Learning Para Clasificacion Con Python}
\item
  \href{https://numerus-scriptum.netlify.app/python/2021-12-13-10-metodos-de-machine-learning-para-regresion-con-python/index.pdf}{\faIcon{file-pdf}}
  \href{https://numerus-scriptum.netlify.app/python/2021-12-13-10-metodos-de-machine-learning-para-regresion-con-python}{10
  Metodos De Machine Learning Para Regresion Con Python}
\item
  \href{https://numerus-scriptum.netlify.app/python/2025-05-10-visualizacion-de-datos-con-python/index.pdf}{\faIcon{file-pdf}}
  \href{https://numerus-scriptum.netlify.app/python/2025-05-10-visualizacion-de-datos-con-python}{Visualizacion
  De Datos Con Python}
\end{enumerate}

Esperamos que encuentres estas publicaciones igualmente interesantes y
útiles. ¡Disfruta de la lectura!






\end{document}
